\section{Related Work}
\label{sec:related}

\paragraph{Specification languages.}
JML~\cite{jml,leavens2006design} provides the contract vocabulary we transport. Eiffel's design-by-contract~\cite{meyer1992design} originated the broader approach; Spec\#~\cite{barnett2005weakest} ported it to .NET. None of these languages address the distribution problem we tackle --- they assume the consumer has source.

\paragraph{Specification inference.}
The inference engine of~\cite{edmonds2026inference} is the immediate predecessor of this work. Daikon~\cite{ernst2007daikon} infers likely invariants from execution traces; Houdini and its descendants infer candidate annotations and discharge them through a verifier. PreInfer~\cite{preinfer2017} uses dynamic symbolic execution. Jdoctor~\cite{jdoctor2018} extracts exception preconditions from Javadoc. None of these systems address what happens to the inferred specifications after they are produced; we provide the missing distribution layer.

\paragraph{Annotation-based contracts.}
JSR-305 captured nullability annotations and was withdrawn before final standardisation; the surviving informal usage covers only nullability. The Checker Framework's annotation vocabulary is type-system-shaped and does not generalise to the relational and quantified clauses JML supports. Cofoja (Contracts for Java) emits runtime checks from \texttt{@Requires}/\texttt{@Ensures} annotations but does not survive without its annotation processor on the classpath; our work is consumer-tolerant by design.

\paragraph{Bytecode metadata transport.}
ASM~\cite{asm} is the standard library for class-file rewriting and underpins our embedder. ByteBuddy and Javassist provide higher-level APIs over similar primitives. None of these libraries prescribe a metadata format for specifications --- they are mechanism, not policy. JEP~181 (nest-mates) and JEP~371 (hidden classes) extend what the JVM accepts as class-file content but do not change the annotation channel.

\paragraph{API description.}
OpenAPI~3.x is the standard for HTTP service descriptions; gRPC's \texttt{.proto} files cover RPC. Pact~\cite{pact} captures consumer-driven contracts as example interactions; it is example-based rather than logical, so it is a transport for examples rather than for specifications, and we specifically reject example-based encoding (Section~\ref{sec:format}). The OpenAPI extension we propose is the smallest viable addition to an existing widely deployed format.

\paragraph{Generating OpenAPI from source.}
Recent work on generating OpenAPI descriptions from Java bytecode and source code includes Respector, Prophet, and AutoOAS~\cite{Lercher2024-AutoOAS}, which extract REST endpoint metadata from Spring Boot applications. These approaches generate the OpenAPI description itself; our \texttt{x-jml-*} extension is orthogonal to that work --- it augments an OpenAPI description with logical contracts, regardless of how the description was produced.

\paragraph{Mutation testing and test generation.}
The mutation-testing literature is cited in detail in~\cite{edmonds2026inference}. The relationship is: that work measured what specifications do for LLM-generated tests, and the present work makes the specifications transportable to LLMs and other consumers that do not have the source.
